% *======================================================================*
%  Cactus Thorn template for ThornGuide documentation
%  Author: Ian Kelley
%  Date: Sun Jun 02, 2002
%  $Header$
%
%  Thorn documentation in the latex file doc/documentation.tex
%  will be included in ThornGuides built with the Cactus make system.
%  The scripts employed by the make system automatically include
%  pages about variables, parameters and scheduling parsed from the
%  relevant thorn CCL files.
%
%  This template contains guidelines which help to assure that your
%  documentation will be correctly added to ThornGuides. More
%  information is available in the Cactus UsersGuide.
%
%  Guidelines:
%   - Do not change anything before the line
%       % START CACTUS THORNGUIDE",
%     except for filling in the title, author, date, etc. fields.
%        - Each of these fields should only be on ONE line.
%        - Author names should be separated with a \\ or a comma.
%   - You can define your own macros, but they must appear after
%     the START CACTUS THORNGUIDE line, and must not redefine standard
%     latex commands.
%   - To avoid name clashes with other thorns, 'labels', 'citations',
%     'references', and 'image' names should conform to the following
%     convention:
%       ARRANGEMENT_THORN_LABEL
%     For example, an image wave.eps in the arrangement CactusWave and
%     thorn WaveToyX should be renamed to CactusWave_WaveToyX_wave.eps
%   - Graphics should only be included using the graphicx package.
%     More specifically, with the "\includegraphics" command.  Do
%     not specify any graphic file extensions in your .tex file. This
%     will allow us to create a PDF version of the ThornGuide
%     via pdflatex.
%   - References should be included with the latex "\bibitem" command.
%   - Use \begin{abstract}...\end{abstract} instead of \abstract{...}
%   - Do not use \appendix, instead include any appendices you need as
%     standard sections.
%   - For the benefit of our Perl scripts, and for future extensions,
%     please use simple latex.
%
% *======================================================================*

\documentclass{article}

% Use the Cactus ThornGuide style file
\usepackage{../../../../doc/latex/cactus}

\begin{document}

\author{Steven R. Brandt}

\title{Cartoon2DX}

\date{September 2 2026}

\maketitle

% Do not delete next line
% START CACTUS THORNGUIDE

\begin{abstract}
\texttt{Cartoon2DX} is the CarpetX port of \texttt{cactusnumerical/Cartoon2D}.
It evolves an axisymmetric system as a thin-in-$y$ Cartesian slab: values on
the $y$-faces (and optionally $x<0$) are obtained by interpolating a
reference profile in cylindrical radius and rotating tensor components into
place. CarpetX treats this as a physical boundary type
\texttt{boundary\_t::cartoon}.
\end{abstract}

\section{Introduction}

Axisymmetric systems can be evolved with a fully 3-dimensional Cartesian
code by keeping only a slab one cell thick in $y$, plus ghost zones, and
filling the off-plane points from the $y=0$ data under the axisymmetry
assumption. This is the ``cartoon'' method of
Alcubierre et al.~\cite{SpacetimeX_Cartoon2DX_Alcubierre2001}, originally
implemented for Cactus3 by Steve Brandt and for Cactus4
(\texttt{Cartoon2D}) by Sai Iyer.

\texttt{Cartoon2DX} does the same thing for the CarpetX driver. The
important CarpetX-specific points are mixed centering (vertex-centered
metric, cell-centered hydro, staggered fluxes) and the fact that
cartoon faces are physical outer faces in CarpetX's
\texttt{valid\_int}/\texttt{valid\_outer}/\texttt{valid\_ghosts}
accounting, not MPI ghosts.

\section{Algorithm}

The symmetry axis is $z$. For a point with coordinates $(x,y)$ (taken
from \texttt{PointDesc.x}/\texttt{.y}, which already know the group's
centering),
\[
\rho = \sqrt{x^2+y^2}.
\]
The reference data live on the $y=0$ plane at $x=\rho$. Each tensor
component is interpolated there with 1D Lagrange interpolation of order
\texttt{Cartoon2DX::order} (using \texttt{order+1} points). The
interpolated Cartesian components are then rotated in the $xy$-plane
into the point's own azimuthal frame:
\begin{itemize}
\item Scalar: copy.
\item Vector $U$: rotate the $x$ and $y$ components; $z$ is unchanged.
\item Symmetric tensor $DD_{\mathrm{sym}}$: the usual quadratic-form
  transform in $xy$; $zz$ is unchanged.
\end{itemize}
The same fill is used for $y$-ghosts and, when
\texttt{fill\_negative\_x=yes}, for $x<0$ hang-over cells.

\section{Using this thorn}

Activate \texttt{Cartoon2DX} and set the cartoon faces on CarpetX:
\begin{verbatim}
CarpetX::boundary_y       = "cartoon"
CarpetX::boundary_upper_y = "cartoon"
CarpetX::boundary_x       = "cartoon"   # half-x; xmin is the axis
CarpetX::ncells_y         = 1
Cartoon2DX::fill_negative_x = yes
Cartoon2DX::order         = 4
\end{verbatim}
\texttt{ghost\_size} must cover the interpolation stencil
(\texttt{order=4} needs five points, so \texttt{ghost\_size $\ge$ 2}).
The one-cell $y$-slab should be centred on $y=0$ (for spacing $dy$,
\texttt{ymin=$-dy/2$}, \texttt{ymax=$+dy/2$}).

CarpetX calls \texttt{ApplyCartoonBoundary(cctkGH, gi, tl)} from
\texttt{apply\_boundary\_conditions} for every live timelevel of a group
that has a cartoon face. Builtin Dirichlet/Neumann/extrapolation
stencils skip those faces. There is no scheduled fill bin and no
\texttt{INHERITS:}.

\subsection{Tensor type}

Each grid-function group is classified from, in order:
\begin{enumerate}
\item the Einstein Toolkit \texttt{tensortypealias} tag
  (\texttt{Scalar}, \texttt{U} or \texttt{D}, \texttt{DD\_sym}) ---
  the same tag old \texttt{Cartoon2D} reads;
\item the CarpetX \texttt{parities} tag (1 / 3 / 6 variables), used by
  generated thorns;
\item untagged one-variable groups, treated as scalars.
\end{enumerate}
Packed groups (for example AsterX \texttt{cons\_vector}) are walked
component-wise from the parities array. Groups with neither tag are
skipped.

\subsection{Half-$x$ and AMR}

Axisymmetry also makes $x<0$ redundant. Set \texttt{xmin=0} and
\texttt{boundary\_x="cartoon"} so the lower-$x$ ghosts are the hang-over
needed for $x$-derivatives; they are filled by the same rotation as the
$y$-ghosts. Do not use \texttt{symmetry\_t::reflection} on $x$ --- that
mirrors Cartesian components, whereas cartoon rotates them.

On AMR, the cartoon slab must stay one cell thick at every level. Set
\texttt{CarpetX::refinement\_factor\_y = 1} and refine only in $x$ and
$z$.

\subsection{Obtaining this thorn}

\texttt{Cartoon2DX} is in the \texttt{SpacetimeX} arrangement. A
GetComponents list that checks out this feature branch together with
the matching CarpetX and AsterX changes is
\texttt{Docs/thornlist/cartoon2dx.th}.

\section{Tests}

\begin{description}
\item[\texttt{test/tov.par}] Cactus testsuite. Low-resolution Cowling
  TOV star, two RK4 steps, intended to finish on a laptop in well under
  30 seconds.
\item[\texttt{test/standalone\_math\_test.cxx}] Interpolation and
  tensor-rotation math, compiled with \texttt{g++}, no Cactus or AMReX.
\end{description}
Longer comparison parameter files (TOV and Brill--Lindquist head-on,
cartoon versus full 3D) are also in \texttt{test/} but are not part of
the testsuite.

\section{History}

\begin{itemize}
\item Cartoon method: Steve Brandt (Cactus3), with Bernd Br\"ugmann.
\item \texttt{Cartoon2D} for Cactus4/Carpet: Sai Iyer, documented by
  Denis Pollney.
\item \texttt{Cartoon2DX} for CarpetX: Steven R. Brandt.
\end{itemize}

An earlier draft of this thorn used \texttt{INHERITS:} of a specific
evolution thorn. That made Cartoon2DX unusable with any other physics
(Cactus treats \texttt{INHERITS:} as a hard activation requirement).
The fill looks up groups by index and writes through
\texttt{CCTK\_VarDataPtrI}, so there is no \texttt{INHERITS:} now.

ENO interpolation and excision, present in old \texttt{Cartoon2D}, are
not implemented.

\begin{thebibliography}{9}

\bibitem{SpacetimeX_Cartoon2DX_Alcubierre2001}
  M.~Alcubierre, S.~Brandt, B.~Br\"ugmann, D.~Holz, E.~Seidel,
  R.~Takahashi and J.~Thornburg,
  ``Symmetry without symmetry: Numerical simulation of axisymmetric
  systems using Cartesian grids,''
  Int.\ J.\ Mod.\ Phys.\ D {\bf 10}, 273 (2001)
  [gr-qc/9908012].

\end{thebibliography}

% Do not delete next line
% END CACTUS THORNGUIDE

\end{document}
